\subsection{Properties of the image of a positive measure}

PROPOSITION 4. --- Let T, T', T'' be three locally compact spaces, μ a positive measure on T, π a μ-measurable mapping of T into T', π' a mapping of T' into T'', and π'' = π' ∘ π.

\begin{enumerate}
\def\labelenumi{\alph{enumi})}
\item
  Suppose that \(\pi\) is \(\mu\) -proper and let \(\mu' = \pi(\mu)\) . For \(\pi'\) to be \(\mu'\) -proper, it is necessary and sufficient that \(\pi''\) be \(\mu\) -proper, in which case \(\pi''(\mu) = \pi'(\pi(\mu))\) (`transitivity of the image of a measure').
\item
  Suppose that \(\pi'\) is continuous, and that \(\pi''\) is \(\mu\) -proper; then \(\pi\) is \(\mu\) -proper, \(\pi'\) is \(\pi(\mu)\) -proper and \(\pi''(\mu) = \pi'(\pi(\mu))\) .
\end{enumerate}

Under the hypotheses of a), for \(\pi''\) to be \(\mu\) -measurable it is necessary and sufficient that \(\pi'\) be \(\mu'\) -measurable (No.~2, Prop. 3). On the other hand if K is a compact subset of \(T''\) , then \(\pi''(K) = \pi^{-1}(\pi'(K))\) ; for \(\pi''(K)\) to be essentially \(\mu\) -integrable, it is necessary and sufficient that \(\pi'(K)\) be essentially \(\mu'\) -integrable, by the Cor. of Th. 1. Finally, if \(\pi''\) is \(\mu\) -proper then, setting \(\mu'' = \pi''(\mu)\) , we have, for every function \(f \in \mathcal{K}(T'')\) ,

\[
\begin{array}{r l} \int f (t ^ {\prime \prime}) d \mu^ {\prime \prime} (t ^ {\prime \prime}) & = \int f (\pi^ {\prime \prime} (t)) d \mu (t) \\ & = \int f (\pi^ {\prime} (\pi (t))) d \mu (t) = \int f (\pi^ {\prime} (t ^ {\prime})) d \mu^ {\prime} (t ^ {\prime}) \end{array}
\]

by Th. 1 of No.~2, which completes the proof of a).

Under the hypotheses of b), let \(K'\) be a compact subset of \(T'\) . Then \(K'' = \pi'(K')\) is compact, therefore \(\overline{\pi''}(K'')\) is essentially \(\mu\) -integrable, therefore \(\overline{\pi}^{1}(K') \subset \overline{\pi''}(K'')\) is essentially \(\mu\) -integrable (Ch. IV, §5, No.~5, Prop. 7), so that \(\pi\) is \(\mu\) -proper. The proof is then concluded by applying part a) of the statement.

COROLLARY. --- Let T and T' be two locally compact spaces, \(\mu\) a positive measure on T, \(\pi\) a bijective mapping of T onto T', and \(\pi^{-1}\) the inverse mapping. Suppose that \(\pi\) is \(\mu\) -proper, and let \(\mu' = \pi(\mu)\) . Then \(\pi^{-1}\) is \(\mu'\) -proper and \(\pi^{-1}(\pi(\mu)) = \mu\) .

PROPOSITION 5. --- Let T and X be two locally compact spaces, \(\mu\) a positive measure on T, \(\pi\) a \(\mu\) -proper mapping of T into X, g a finite numerical function \(\geqslant 0\) , defined on X and such that \(g \circ \pi\) is locally integrable for \(\mu\) . In order that g be locally integrable for \(\pi(\mu)\) , it is necessary and sufficient that \(\pi\) be proper for the measure \((g \circ \pi) \cdot \mu\) , in which case

\[
\pi \big ((g \circ \pi) \cdot \mu \big) = g \cdot \pi (\mu).\tag{4}
\]

Set \(\nu = \pi(\mu)\) . For \(g\) to be locally \(\nu\) -integrable, it is necessary and sufficient that \(gf\) be \(\nu\) -integrable for every function \(f \in \mathcal{K}(X)\) ; since \(gf\) has compact support, it comes to the same to say that \(gf\) is essentially \(\nu\) -integrable, and this is equivalent to saying that \((g \circ \pi)(f \circ \pi)\) is essentially \(\mu\) -integrable (Th. 1). But, by Th. 1 of §5, No.~3, this signifies that \(f \circ \pi\) is essentially integrable for \(\rho = (g \circ \pi) \cdot \mu\) , and, by definition, this says that \(\pi\) is \(\rho\) -proper (since \(\pi\) is obviously \(\rho\) -measurable). Moreover,

\[
\int f g d \nu = \int f (\pi (t)) g (\pi (t)) d \mu (t) = \int f (\pi (t)) d \rho (t)
\]

(No.~2, Th. 1 and §5, Th. 1), which proves the relation (4).

PROPOSITION 6. --- Let T and X be two locally compact spaces, \((\lambda_{\alpha})_{\alpha\in\mathbb{A}}\) a family of positive measures on T, directed for the relation \(\leqslant\) , admitting in \(\mathcal{M}(\mathrm{T})\) a supremum \(\mu\) . In order that a mapping \(\pi\) of T into X be \(\mu\) -proper, it is necessary and sufficient that it be \(\lambda_{\alpha}\) -proper for all \(\alpha\in A\) , and that the family \(\left(\pi(\lambda_{\alpha})\right)_{\alpha\in\mathbb{A}}\) be bounded above in \(\mathcal{M}(\mathrm{X})\) . In this case,

\[
\pi (\mu) = \sup _ {\alpha} \pi (\lambda_ {\alpha}).\tag{5}
\]

For \(\pi\) to be \(\mu\) -measurable, it is necessary and sufficient that \(\pi\) be \(\lambda_{\alpha}\) -measurable for all \(\alpha \in A\) (§1, No.~4, Cor. 2 of Prop. 11). Suppose that this condition is satisfied; to say that \(\pi\) is \(\mu\) -proper is then equivalent to saying that, for every function \(f \in \mathcal{K}_{+}(X)\) ,

\[
\mu^ {\bullet} (f \circ \pi) <   + \infty .
\]

Now,

\[
\int^ {\bullet} (f \circ \pi) d \mu = \sup _ {\alpha} \int^ {\bullet} (f \circ \pi) d \lambda_ {\alpha} = \sup _ {\alpha} \int^ {\bullet} f d (\pi (\lambda_ {\alpha}))
\]

(§1, No.~4, Prop. 11); the first member is thus finite for every \(f \in \mathcal{K}_{+}(X)\) if and only if the family \((\pi(\lambda_{\alpha}))\) admits a supremum \(\theta\) in \(\mathcal{M}(X)\) , in which case \(\int(f \circ \pi) d\mu = \int f d\theta\) , a relation equivalent to (5).

COROLLARY 1. --- Let \((\mu_{\alpha})_{\alpha \in A}\) be a summable family of positive measures on T, such that \(\mu = \sum_{\alpha \in A} \mu_{\alpha}\) ; for a mapping \(\pi\) of T into a locally compact space X to be \(\mu\) -proper, it is necessary and sufficient that it be \(\mu_{\alpha}\) -proper for every \(\alpha \in A\) , and that the family \((\pi(\mu_{\alpha}))_{\alpha \in A}\) be summable. In this case,

\[
\pi (\mu) = \sum_ {\alpha \in A} \pi (\mu_ {\alpha}).\tag{6}
\]

COROLLARY 2. --- Let T and X be two locally compact spaces, \((\lambda_{i})_{1\leqslant i\leqslant n}\) a finite sequence of positive measures on T, and let \(\mu=\sum_{i=1}^{n}\lambda_{i}\) . For a mapping \(\pi\) of T into X to be \(\mu\) -proper, it is necessary and sufficient that it be \(\lambda_{i}\) -proper for every index i, in which case

\[
\sum_ {i = 1} ^ {n} \pi (\lambda_ {i}) = \pi \left(\sum_ {i = 1} ^ {n} \lambda_ {i}\right).
\]

